% titel: Entleerung eines Behälters
% skript: Kap. Hydrodynamik

% Ein Behälter mit Öl entleert sich über ein horizontales Rohr am Boden;
% gefragt wird nach Ausflussgeschwindigkeit und Füllhöhe, zunächst ideal
% reibungsfrei (Torricelli), dann mit laminarer Reibung im Rohr
% (Hagen-Poiseuille).

Ein zylindrischer Behälter (Radius $r_1$) entleert sich über ein waagrechtes
Rohr (Radius $r_2\ll r_1$, Länge $l$) am Boden. Gesucht werden die
Ausflussgeschwindigkeit und der zeitliche Verlauf der Füllhöhe $z(t)$ ---
zuerst für reibungsfreies, dann für laminar reibend strömendes Öl.

\begin{werkzeugkasten}[Werkzeuge]
\textbf{Torricelli / Bernoulli:} Für ein offenes Gefäss (oberhalb der
Oberfläche und an der Ausflussöffnung herrscht derselbe Umgebungsdruck) und
$r_2\ll r_1$ (Oberfläche ruht, $v_1\approx 0$) folgt die
Ausflussgeschwindigkeit aus Bernoulli:
\begin{equation}
  v_2 \approx \sqrt{2gz}\;,
\end{equation}
mit der momentanen Füllhöhe $z$ (Höhe der Säule über der Ausflussöffnung).

\textbf{Kontinuität} (Volumenstromerhalt), auf einen sich leerenden Tank
angewandt ($\dot z\equiv \mathrm dz/\mathrm dt$):
\begin{equation}
  A_1\,\dot z = -A_2\,v_2\;,\qquad
  A_1=\pi r_1^2,\; A_2=\pi r_2^2\;.
\end{equation}

\textbf{Hagen-Poiseuille} (laminare, stationäre Rohrströmung): Volumenstrom
und mittlere Geschwindigkeit durch ein Rohr der Länge $l$ mit
Druckdifferenz $\Delta p$
\begin{equation}
  \dot V=\frac{\pi\,r_2^4}{8\,\eta}\,\frac{\Delta p}{l}\;,\qquad
  \bar u=\frac{\dot V}{A_2}
  =\frac{r_2^2}{8\,\eta}\,\frac{\Delta p}{l}\;,
\end{equation}
mit der dynamischen Viskosität $\eta$.

\textbf{Einheiten und Bezeichner:} $\dot z$ in \unit{\metre\per\second},
$A$ in \unit{\square\metre}, $\eta$ in \unit{\pascal\second},
$\Delta p$ in \unit{\pascal}.
\end{werkzeugkasten}

\begin{werkzeugkasten}[Vorgehen in allgemeiner Form]
\begin{itemize}
  \item \emph{Ideal (a--c):} $v_2$ aus Torricelli einsetzen, Kontinuität liefert
       eine gewöhnliche DGL erster Ordnung für $z(t)$; mit Trennung der
       Veränderlichen integrieren, Anfangsbedingung $z(0)=z_0$.
  \item \emph{Reibend (d--e):} $v_2=\bar u$ aus Hagen-Poiseuille; der
       quadratische (dynamische) Druckterm ist gegenüber dem linearen zu
       vernachlässigen, sodass eine lineare DGL entsteht.
\end{itemize}
\end{werkzeugkasten}

\begin{achtung}[Typische Fallen]
\begin{itemize}
  \item Vorzeichen in der Kontinuität: wegen $\dot z<0$ wird der Tank leer.
  \item Bei Trennung der Veränderlichen die Integrationsgrenzen der Umdrehung
       ($z(t_0)$ nach $z(t_1)$) korrekt setzen --- hier als Hinweis gegeben.
  \item In Teil (d) ist $\bar u$ eine \emph{Anfangs}-Grösse; sie hängt von $z_0$
       ab, nicht mit der Torricelli-Geschwindigkeit verwechseln.
  \item Viskosität umrechnen: \unit{\milli\pascal\second}
       $\to$ \unit{\pascal\second}.
\end{itemize}
\end{achtung}
